%% ma: L12-B06-0021
%% lop: 12
%% chuong: 2
%% bai: 6
%% dang: 12.6.14
%% loai: TN4
%% muc: TH
%% dap_an: "C"
%% nguon: "LaTeX giáo viên gửi 21/10/2026 (M05, Vectơ và hệ trục tọa độ trong không gian, 107 câu), Phần I Câu 29"
%% kien_thuc: "Bài 6 (tổng vectơ, bình phương vô hướng)"
%% trang_thai: chinh-thuc
%% ghi_chu: "Giáo viên đã duyệt 21/10/2026. Phần I Câu 45 của tài liệu (ba lực 6 N, 8 N, 24 N) chỉ đổi số so với câu này nên đã bỏ"
\begin{ex}
Một vật chịu tác dụng đồng thời của $4$ lực $\overrightarrow{F_1}$, $\overrightarrow{F_2}$, $\overrightarrow{F_3}$, $\overrightarrow{F_4}$ và ở trạng thái cân bằng. Biết rằng $\overrightarrow{F_1}$, $\overrightarrow{F_2}$, $\overrightarrow{F_3}$ là ba vectơ lực đôi một vuông góc và có cùng độ lớn là $10$ N. Độ lớn của lực $\overrightarrow{F_4}$ bằng
\choice{$10$ N}{$10\sqrt2$ N}{\True $10\sqrt3$ N}{$30$ N}
\loigiai{Vật cân bằng nên $\overrightarrow{F_4}=-(\overrightarrow{F_1}+\overrightarrow{F_2}+\overrightarrow{F_3})$. Ba lực đôi một vuông góc nên $|\overrightarrow{F_1}+\overrightarrow{F_2}+\overrightarrow{F_3}|^2=3\cdot10^2$. Vậy $|\overrightarrow{F_4}|=10\sqrt3$ N.}
\end{ex}
